\section{Exact rational padding and certified signing}
\label{sec:signing}


The deterministic signing theorem proved in Appendix~\ref{app:signing}
has an exact input promise. For explicit
rational columns $v_1,\ldots,v_M\in(\Q+i\Q)^d$ satisfying
$\sum_i v_iv_i^*=I$, it returns signs $\sigma_i\in\{-1,1\}$ with
\begin{equation}
 \left\|\sum_i\sigma_i v_iv_i^*\right\|_{\op}
       \le40\sqrt{\Delta},\qquad
       \Delta=\max_i\|v_i\|^2,
 \label{eq:ta-signing-input}
\end{equation}
in time polynomial in the full binary input length. This is
\Cref{thm:rational-signing}, a quantitative form of the method of
Ezeunala and Jiang~\cite{EzeunalaJiang2026}. Real rational columns are a
special case. We first establish an elementary exact factorization, then prove the
transfer required here, including identity completion and a rational
success certificate.

\begin{lemma}[Rational positive semidefinite decomposition]
\label{lem:rational-psd}
Every rational Hermitian positive semidefinite matrix $A$ can be written
exactly as
\[
                         A=\sum_{j=1}^M v_jv_j^*,
 \qquad v_j\in\Q(\mathrm i)^n,
\]
in polynomial bit time, with $M$ and the total encoding length of the
columns polynomial in the encoding length of $A$.  If $A$ is real,
the columns may be taken real.  For $A=0$, the sum may be empty.
\end{lemma}

\begin{proof}
If a positive semidefinite matrix has all diagonal entries zero, it is
zero, since its two-by-two principal minors give
$|a_{ij}|^2\le a_{ii}a_{jj}$.  Otherwise permute a positive diagonal
entry into the first position and write
\[
 A=\begin{pmatrix}d&b^*\\ b&C\end{pmatrix}
  =d\begin{pmatrix}1\\b/d\end{pmatrix}
       \begin{pmatrix}1\\b/d\end{pmatrix}^{\!*}
       +\begin{pmatrix}0&0\\0&C-bb^*/d\end{pmatrix}.
\]
Here $d$ is positive rational and the Schur complement is positive
semidefinite: evaluate the quadratic form of $A$ at
$(-b^*x/d,x)$.  Repetition gives a rational pivoted $LDL^*$
factorization, or equivalently $A=\sum_j d_j u_ju_j^*$ with positive
rational $d_j$ and rational complex $u_j$.

All elimination entries have polynomial encoding length.  To see this
directly, clear the input denominators by one positive integer $q$.
The resulting Gaussian-integer matrix has polynomial-bit entries,
and every minor has polynomial bit length by the determinant
expansion.  At any pivot order, the pivots and elimination coefficients
are ratios of these minors, with an additional factor $q$ for the
pivots.  There are at most $n$ pivots, so rational elimination takes
polynomial bit time.

For a pivot $d=a/b>0$ with positive integers $a,b$, expand
$ab=\sum_k c_k4^k$, where $c_k\in\{0,1,2,3\}$.  Then
\[
                 d=\sum_k c_k\left(\frac{2^k}{b}\right)^2.
\]
Replace $d u u^*$ by $c_k$ copies of the outer product of
$(2^k/b)u$ for each $k$.  The number of copies and their bit lengths
are polynomial in those of $a,b,u$.  This proves every assertion,
including the real case.
\end{proof}


\begin{lemma}[Checked signing of the stacked frame]
\label{lem:ta-signing}
Let $X$ and the union $B$ of $q$ trees be as in
\Cref{prop:ta-extraction}, retaining $r/2\le q\le r$ trees, and let
$z_e$ be the frame of \Cref{lem:ta-frame}. Put
\begin{equation}
                   w_e=(X^{-1/2}b_e,z_e)\qquad(e\in B).
 \label{eq:ta-stacked}
\end{equation}
For an absolute constant $C_s$, a deterministic polynomial-time procedure
finds $A\subseteq B$ satisfying
\begin{equation}
 \left\|\sum_{e\in A}w_ew_e^T
                    -\frac12\sum_{e\in B}w_ew_e^T\right\|_{\op}
                          \le C_s r^{-1/8}.
 \label{eq:ta-half-frame}
\end{equation}
An exact polynomial-time rational PSD test certifies the output.
\end{lemma}

\begin{proof}
The first block of \Cref{eq:ta-stacked} is expressed in orthonormal
coordinates on $\one^\perp$; its particular orthogonal basis is
irrelevant. Since $L_B\preceq L_H\preceq X$, its outer-product sum is
at most $I$, while the second block sums to $I$. Block
Cauchy--Schwarz therefore gives
\begin{equation}
                 \sum_{e\in B}w_ew_e^T\preceq2I.
 \label{eq:ta-subisotropic}
\end{equation}
Individual squared norms are at most
$r^{-1/4}+(1+4/r)/q\le5r^{-1/4}$ for the chosen cutoff.

Write $n=|V(H)|$, $m=|E(H)|$, and $P=n+m+2$, as in
\Cref{lem:ta-bits}. The centered positive definite $X$ has polynomial-bit rational entries,
with smallest eigenvalue at least $P^{-3}$ and norm at most $6m$.
As in the proof of \Cref{lem:ta-frame}, elimination, square-root
approximation, and solves compute a whitening of this block to any
requested polynomial number of bits. More explicitly, in the rational
coordinates used above write $X=CA_XC^T$ and $b_e=Cc_e$. Then
$c_e^TA_X^{-1}c_f=b_e^TX^{-1}b_f$, so whitening $A_X$ gives the same
Gram matrix as the displayed block and differs from it only by an
orthogonal map. The smallest eigenvalue of $A_X$ is at least
$P^{-3}/n$, because the centered projection of $(v,0)$ has squared norm
at least $\|v\|^2/n$. Thus this computation has the same polynomial
precision bounds. Combining it with \Cref{lem:ta-frame}, compute
rational columns $a_e$ approximating $w_e/2$ to Euclidean error at most
$1/(1000|B|r)$. Because \Cref{eq:ta-subisotropic} implies
$\|w_e/2\|\le1$, the outer-product identity
\[
 aa^T-vv^T=(a-v)a^T+v(a-v)^T
\]
shows that
\begin{equation}
 \sum_{e\in B}\|a_ea_e^T-w_ew_e^T/4\|_{\op}
       \le\frac1{100r},\qquad
 \Sigma:=\sum_{e\in B}a_ea_e^T\preceq\frac34I.
 \label{eq:ta-approximate-data}
\end{equation}
The same estimate controls every signed sum, without knowing the signs.
Also $\|a_e\|^2\le3r^{-1/4}$.

Apply \Cref{lem:rational-psd} to the rational PSD residual
$I-\Sigma$, obtaining polynomially many polynomial-bit real rational
columns $v_1,\ldots,v_J$ with
\begin{equation}
                      VV^T=\sum_{j=1}^Jv_jv_j^T=I-\Sigma,
                         \qquad V=[v_1\ \cdots\ v_J].
 \label{eq:ta-padding-factor}
\end{equation}
In particular $\|V\|_{\op}\le1$ and each $\|v_j\|\le1$.
Extend the data columns by zeros in $J$ additional coordinates. Set
$N=2r$. For each $j$, add $N^2$ columns
\begin{equation}
                 p_{j,k}=\frac1N(v_j,\tau_{j,k}f_j),
       \qquad 1\le k\le N^2,
 \label{eq:ta-padding-columns}
\end{equation}
where $f_j$ is the $j$th extra coordinate vector and exactly half the
$\tau_{j,k}$ equal $1$, half $-1$. The parity requirement holds because
$N$ is even. Their unsigned sum is
\[
 \sum_{j,k}p_{j,k}p_{j,k}^T
       =\begin{pmatrix}I-\Sigma&0\\0&I_J\end{pmatrix}.
\]
Together with the extended data this is exactly the identity.
Every padding squared norm is at most $2/N^2$, so the maximum squared
norm $\Delta$ of the entire rational system satisfies
$\Delta\le3r^{-1/4}$. The number $|B|+JN^2$ of explicit columns and
their binary lengths are polynomial, since $r\le m$.

Apply \Cref{eq:ta-signing-input}. Write $S$ for the augmented signed
matrix, and distinguish the output signs $\sigma_{j,k}$ from the fixed
construction signs $\tau_{j,k}$. Its lower principal block is
\[
 D_0=\diag\left(N^{-2}\sum_k\sigma_{1,k},\ldots,
                         N^{-2}\sum_k\sigma_{J,k}\right).
\]
The padding contribution to its upper principal block is exactly
\begin{equation}
                              VD_0V^T.
 \label{eq:ta-padding-removal}
\end{equation}
If $\|S\|\le T$, then $\|D_0\|\le T$ and
$\|VD_0V^T\|\le T$. Subtracting the latter from the full upper block
gives
\begin{equation}
                    \left\|\sum_{e\in B}\sigma_ea_ea_e^T\right\|
                                \le2T.
 \label{eq:ta-data-signing}
\end{equation}
No assertion about the signed off-diagonal block is needed.

For a fully rational certificate, choose a dyadic $d_0$ with
$\sqrt\Delta\le d_0\le2\sqrt\Delta$ by squared comparisons and set
$T=2^{15}d_0$. Check $TI-S\succeq0$ and $TI+S\succeq0$ exactly by
rational elimination. Since $40<2^{15}$, the signing theorem guarantees
that both tests pass. The certified output satisfies
\Cref{eq:ta-data-signing} and $T\le2^{16}\sqrt\Delta$.

Return the positive data signs $A=\{e:\sigma_e=1\}$.
By \Cref{eq:ta-approximate-data,eq:ta-data-signing}, its centered half
sum has norm at most
\[
 \frac12\left\|\sum_e\sigma_e w_ew_e^T\right\|
       \le4T+\frac2{100r}
       \le2^{20}r^{-1/8}.
\]
Thus one may take $C_s=2^{20}$. The rational PSD tests have polynomial
cost and certify all the required matrix estimates simultaneously.
The signing algorithm and all numerical routines above have prescribed
polynomial precision and iteration bounds. This proves deterministic
polynomial bit complexity.
\end{proof}
